% !TEX program = pdflatex
% SHILATE Bachelor's Thesis
% Software-Defined Vehicle Reinforcement Learning: From Simulation to SDV Deployment
% Author: [Candidate Name]
% Supervisor: [Supervisor Name]
% Politehnica University of Timisoara, 2026

% ============================================================
%  PREAMBLE
% ============================================================
\PassOptionsToPackage{unicode}{hyperref}
\PassOptionsToPackage{hyphens}{url}

\documentclass[12pt,a4paper]{article}

% --- Encoding & fonts ---
\usepackage[T1]{fontenc}
\usepackage[utf8]{inputenc}
\usepackage{lmodern}
\usepackage{textcomp}

% --- Page geometry (UPT thesis: mirror margins) ---
\usepackage[
  a4paper,
  top=2.5cm,
  bottom=2cm,
  inner=2cm,
  outer=2cm,
  includeheadfoot
]{geometry}

% --- Line spacing ---
\usepackage{setspace}
\setstretch{1.15}

% --- Headers / footers ---
\usepackage{fancyhdr}
\pagestyle{fancy}
\fancyhf{}
\fancyhead[L]{\fontsize{8}{10}\selectfont
  Politehnica University of Timisoara \\
  Computer Engineering, 2026 \\
  \leftmark}
\fancyfoot[C]{\thepage}
\renewcommand{\headrulewidth}{0.4pt}
\setlength{\headheight}{38pt}

% --- Mathematics ---
\usepackage{amsmath,amssymb}

% --- Graphics ---
\usepackage{graphicx}
\usepackage{float}
\graphicspath{{figures/}}
\makeatletter
\def\maxwidth{\ifdim\Gin@nat@width>\linewidth\linewidth\else\Gin@nat@width\fi}
\def\maxheight{\ifdim\Gin@nat@height>\textheight\textheight\else\Gin@nat@height\fi}
\makeatother
\setkeys{Gin}{width=\maxwidth,height=\maxheight,keepaspectratio}

% --- Tables ---
\usepackage{longtable,booktabs,array,multirow}
\usepackage{calc}

% --- Lists & spacing ---
\usepackage{enumitem}
\setlength{\parindent}{1.5cm}
\setlength{\parskip}{0pt}

% --- Code listings ---
\usepackage{listings}
\usepackage{xcolor}
\definecolor{codegray}{rgb}{0.5,0.5,0.5}
\definecolor{codeblue}{rgb}{0.0,0.3,0.7}
\definecolor{codegreen}{rgb}{0.0,0.5,0.0}
\definecolor{codeback}{rgb}{0.97,0.97,0.97}
\lstset{
  backgroundcolor=\color{codeback},
  basicstyle=\ttfamily\small,
  breakatwhitespace=false,
  breaklines=true,
  captionpos=b,
  commentstyle=\color{codegreen},
  keepspaces=true,
  keywordstyle=\color{codeblue}\bfseries,
  numberstyle=\tiny\color{codegray},
  numbers=left,
  stringstyle=\color{codegray},
  tabsize=2,
  frame=single,
  rulecolor=\color{codegray},
}

% --- Miscellaneous ---
\usepackage{microtype}
\usepackage{etoolbox}
\usepackage{footnotehyper}
\makesavenoteenv{longtable}
\setlength{\emergencystretch}{3em}
\providecommand{\tightlist}{%
  \setlength{\itemsep}{0pt}\setlength{\parskip}{0pt}}

% --- Bibliography ---
\usepackage[
  backend=bibtex,
  style=numeric,
  sorting=nyt,
  citestyle=numeric-comp
]{biblatex}
\addbibresource{shilate_thesis.bib}

% --- Hyperlinks (last) ---
\usepackage{bookmark}
\usepackage{hyperref}
\hypersetup{
  colorlinks=true,
  linkcolor=black,
  citecolor=black,
  urlcolor=black,
  pdftitle={SHILATE: Software-Defined Vehicle Reinforcement Learning},
  pdfauthor={[Candidate Name]},
}
\usepackage{xurl}
\urlstyle{same}

% ============================================================
%  TITLE PAGE DATA
% ============================================================
\title{}
\author{}
\date{}

% ============================================================
%  DOCUMENT START
% ============================================================
\begin{document}

% Suppress page number on title and abstract pages
\pagenumbering{roman}
\pagestyle{empty}

% ============================================================
%  TITLE PAGE
% ============================================================
\begin{titlepage}
  \centering
  \vspace*{2cm}

  {\fontsize{10}{13}\selectfont
    \textbf{POLITEHNICA UNIVERSITY OF TIMISOARA}\\[0.3cm]
    \textbf{Faculty of Automation and Computers}\\[0.3cm]
    \textbf{Computer Engineering Study Programme}
  }

  \vspace{2.5cm}

  {\fontsize{20}{26}\selectfont\bfseries\MakeUppercase{%
    SHILATE: Software-Defined Vehicle\\[0.4cm]
    Reinforcement Learning ---\\[0.4cm]
    From Simulation to SDV Deployment%
  }\par}

  \vspace{1.5cm}

  {\fontsize{14}{18}\selectfont
    \textbf{BACHELOR'S THESIS}
  }

  \vspace{3cm}

  \begin{flushleft}
    {\fontsize{14}{18}\selectfont
      \textbf{Candidate: Firstname, Lastname}\\[0.8cm]
      \textbf{Scientific Supervisor: Prof.dr.ing. Firstname, Lastname}
    }
  \end{flushleft}

  \vfill

  {\fontsize{14}{18}\selectfont June 2026}
\end{titlepage}

\clearpage

% ============================================================
%  ABSTRACT
% ============================================================
\vspace*{3\baselineskip}

\begin{center}
  {\fontsize{20}{26}\selectfont\bfseries ABSTRACT}
\end{center}

\vspace{2\baselineskip}

The proliferation of software-defined vehicle (SDV) platforms has opened an
opportunity to combine modern deep reinforcement learning (RL) with real
automotive runtime stacks. This thesis presents \textbf{SHILATE} --- a
complete, open-source pipeline that trains a Proximal Policy Optimization (PPO)
agent inside a Unity physics simulation and deploys the resulting policy onto an
Eclipse Leda SDV stack, bridging the gap between laboratory RL research and
realistic vehicle software deployment.

The simulation component is built in Unity 2022 LTS. A single rear-wheel-drive
vehicle navigates a circular obstacle course instrumented with a 21-ray
LiDAR-style sensor. The training bridge computes a shaped reward signal from
progress, collision events, and episode milestones, and communicates with the
Python training process exclusively over MQTT, keeping the two subsystems fully
decoupled. A dedicated Editor window (\emph{SHILATE Training Controller})
provides real-time health monitoring, metric graphs, and one-click training
control without leaving the Unity Editor.

The Python training stack wraps the simulation in a Gymnasium-compliant
environment, configures Stable-Baselines3 PPO with a two-layer actor-critic
network, and streams structured telemetry back to the Editor. Trained
checkpoints are subsequently deployed to the Eclipse Leda SDV image via a
containerised MQTT-to-Kuksa bridge and a Velocitas vehicle application that
enforces operational limits (overspeed, RPM redline) on the running agent.

Experimental results demonstrate that the PPO agent converges reliably on the
obstacle course within 100k interaction steps and that the end-to-end deployment
chain correctly propagates vehicle telemetry from the simulation through the
VSS data layer to the Velocitas safety monitor. The project shows that an
integrated sim-to-SDV pipeline is achievable with commodity open-source
tooling, and identifies curriculum learning and multi-environment
parallelisation as natural extensions for future work.

\textbf{Keywords:} reinforcement learning, software-defined vehicle, Unity
simulation, PPO, MQTT, Eclipse Leda, Kuksa, Stable-Baselines3, Gymnasium.

\clearpage

% ============================================================
%  TABLE OF CONTENTS
% ============================================================
\pagestyle{fancy}
\tableofcontents
\clearpage

% Switch to Arabic page numbering starting from the Introduction
\pagenumbering{arabic}
\setcounter{page}{1}

% ============================================================
%  CHAPTER 1 — INTRODUCTION
% ============================================================
\markboth{Chapter 1 — Introduction}{}
\section*{\centering CHAPTER 1\\[0.3cm] INTRODUCTION}
\addcontentsline{toc}{section}{Chapter 1 \quad Introduction}
\vspace{2\baselineskip}

\subsection*{1.1 \quad MOTIVATION AND CONTEXT}
\addcontentsline{toc}{subsection}{1.1 \quad Motivation and Context}

\indent The automotive industry is undergoing a fundamental architectural shift.
Traditional vehicles are increasingly giving way to software-defined vehicles
(SDVs) --- platforms in which electronic control units are consolidated onto
high-performance compute hardware and the vehicle's behaviour is governed by
software layers that can be updated over the air. This transformation enables
capabilities that were previously inconceivable: dynamic reconfiguration of
driving characteristics, continuous safety-critical updates without a workshop
visit, and the integration of machine-learning-based components directly into
the vehicle runtime.

At the same time, deep reinforcement learning (RL) has matured into a practical
engineering tool. Algorithms such as Proximal Policy Optimization (PPO)
\cite{Schulman2017} achieve human-competitive or superhuman performance in
continuous control tasks ranging from robotic locomotion to game playing, with
comparatively modest tuning effort. The natural question that arises is
therefore: \emph{how can an RL-trained policy be moved, end-to-end, from a
training simulation onto an SDV runtime stack in a reproducible and maintainable
way?}

This question is not merely academic. Automotive-grade deployment requires that
a trained model run inside a standardised software container, interact with
vehicle signals through a well-defined data layer, and remain subject to
functional-safety monitors that can intervene when the model operates outside
its intended envelope. Bridging this gap between laboratory RL and production
SDV middleware is the core motivation behind the work presented in this thesis.

\subsection*{1.2 \quad PROBLEM STATEMENT}
\addcontentsline{toc}{subsection}{1.2 \quad Problem Statement}

\indent Existing RL research in the autonomous driving domain concentrates
almost exclusively on training performance --- reward curves, sample efficiency,
and policy robustness. The software engineering challenge of taking a trained
checkpoint and deploying it as a first-class citizen of an SDV runtime stack is
largely unexplored in the open-source community. Concretely, the following
sub-problems remain unaddressed in an integrated, open-source pipeline:

\begin{enumerate}[leftmargin=2.5cm, label=\textbf{P\arabic*.}]
  \item \textbf{Simulation-to-middleware coupling.} A physics simulator and an
    SDV middleware stack speak fundamentally different protocols; a
    communication bridge that is both low-latency and decoupled from either
    system is needed.
  \item \textbf{Standardised vehicle-signal representation.} Raw simulator
    outputs (floating-point physics quantities) must be mapped onto
    industry-standard vehicle signal specifications (VSS) so that the SDV
    safety layer can interpret them without simulator-specific knowledge.
  \item \textbf{Operational safety monitoring.} A deployed RL policy may
    extrapolate outside its training distribution; a lightweight, declarative
    safety monitor must be able to override the policy's actions in real time.
  \item \textbf{Reproducible training workflow.} The entire pipeline --- from
    launching the simulation to saving a deployable checkpoint --- must be
    reproducible from a single control surface without manual orchestration.
\end{enumerate}

\subsection*{1.3 \quad OBJECTIVES}
\addcontentsline{toc}{subsection}{1.3 \quad Objectives}

\indent This thesis addresses the above problems through the design, implementation, and
evaluation of \textbf{SHILATE} (Software-defined-vehicle Hardware-In-the-Loop
Autonomous Training Environment). The specific objectives are:

\begin{enumerate}[leftmargin=2.5cm, label=\textbf{O\arabic*.}]
  \item Design and implement a Unity-based physics simulation of a rear-wheel-drive
    vehicle navigating an instrumented obstacle course.
  \item Define a shaped reward function that guides a PPO agent to complete
    laps reliably and safely, and expose a Gymnasium-compliant Python interface
    over MQTT.
  \item Implement a real-time health-monitoring and training-control layer inside
    the Unity Editor so that a researcher can observe and intervene in a training
    run without leaving the development environment.
  \item Construct an SDV deployment chain using Eclipse Leda, the Kuksa vehicle
    signal databroker, and a Velocitas vehicle application that enforces
    operational limits on the running policy.
  \item Evaluate the pipeline quantitatively: training convergence, inference
    lap performance, and end-to-end signal propagation through the SDV stack.
\end{enumerate}

\subsection*{1.4 \quad THESIS STRUCTURE}
\addcontentsline{toc}{subsection}{1.4 \quad Thesis Structure}

\indent The remainder of this document is organised as follows. Chapter~2 reviews the
theoretical background in reinforcement learning, SDV middleware, and
communication protocols, and situates SHILATE among related work. Chapter~3
presents the overall system architecture and the rationale for key design
decisions. Chapter~4 describes the Unity simulation in detail, covering the
vehicle model, perception layer, reward function, and training control
interface. Chapter~5 documents the Python RL training pipeline, including the
Gymnasium wrapper, PPO configuration, and health telemetry. Chapter~6 explains
the SDV integration through Eclipse Leda, focusing on signal mapping and
safety monitoring. Chapter~7 reports experimental results. Chapter~8 concludes
with a summary of contributions, limitations, and directions for future work.

\clearpage

% ============================================================
%  CHAPTER 2 — BACKGROUND AND RELATED WORK
% ============================================================
\markboth{Chapter 2 — Background \& Related Work}{}
\section*{\centering CHAPTER 2\\[0.3cm] BACKGROUND AND RELATED WORK}
\addcontentsline{toc}{section}{Chapter 2 \quad Background and Related Work}
\vspace{2\baselineskip}

\subsection*{2.1 \quad REINFORCEMENT LEARNING FUNDAMENTALS}
\addcontentsline{toc}{subsection}{2.1 \quad Reinforcement Learning Fundamentals}

\indent Reinforcement learning is a computational framework for sequential decision
making under uncertainty \cite{Sutton2018}. An agent interacts with an
environment modelled as a Markov Decision Process (MDP), defined by the tuple
$(\mathcal{S}, \mathcal{A}, P, R, \gamma)$, where $\mathcal{S}$ is the state
space, $\mathcal{A}$ is the action space, $P: \mathcal{S} \times \mathcal{A}
\times \mathcal{S} \rightarrow [0,1]$ is the transition kernel, $R:
\mathcal{S} \times \mathcal{A} \rightarrow \mathbb{R}$ is the reward function,
and $\gamma \in [0,1)$ is the discount factor. The agent's goal is to learn a
policy $\pi: \mathcal{S} \rightarrow \mathcal{A}$ (or a distribution over
actions) that maximises the expected discounted return $G_t = \sum_{k=0}^{\infty}
\gamma^k R_{t+k+1}$.

\textbf{Policy gradient methods} optimise the policy directly by computing the
gradient of the expected return with respect to the policy parameters
$\theta$. The foundational result is the policy gradient theorem:

\begin{equation}
  \nabla_\theta J(\theta) = \mathbb{E}_{\pi_\theta}
    \left[ \nabla_\theta \log \pi_\theta(a|s) \cdot Q^{\pi_\theta}(s, a) \right],
\end{equation}

\noindent where $Q^{\pi_\theta}(s,a)$ is the action-value function under the
current policy. In practice, a learned value function (critic) $V_\phi(s)$
is used to compute the advantage $A(s,a) = Q(s,a) - V_\phi(s)$, which
reduces gradient variance without introducing bias.

\textbf{Actor-critic} architectures maintain both a policy network (actor) and a
value network (critic), trained simultaneously. This combination underpins
most modern continuous-control RL algorithms, including the one used in this
work.

\subsection*{2.2 \quad PROXIMAL POLICY OPTIMIZATION}
\addcontentsline{toc}{subsection}{2.2 \quad Proximal Policy Optimization}

\indent Proximal Policy Optimization \cite{Schulman2017} is a policy gradient algorithm
designed to take the largest possible improvement step on a policy without
destabilising training. It addresses the instability of earlier methods such
as TRPO by replacing the constrained optimisation with a clipped surrogate
objective that is cheap to compute and easy to implement:

\begin{equation}
  L^{\text{CLIP}}(\theta) = \mathbb{E}_t \left[
    \min\!\left( r_t(\theta)\,\hat{A}_t,\;
      \text{clip}\!\left(r_t(\theta),\,1-\epsilon,\,1+\epsilon\right)\hat{A}_t
    \right)
  \right],
\end{equation}

\noindent where $r_t(\theta) = \pi_\theta(a_t|s_t) / \pi_{\theta_\text{old}}(a_t|s_t)$
is the probability ratio between the updated and old policies, $\hat{A}_t$ is
the estimated advantage at time step $t$, and $\epsilon$ is a clipping
hyperparameter (typically $0.2$).

The full training objective combines the clipped policy loss, a value function
regression term, and an entropy bonus:
\begin{equation}
  L(\theta) = L^{\text{CLIP}}(\theta)
            - c_1 L^{\text{VF}}(\theta)
            + c_2 H[\pi_\theta(\cdot|s_t)],
\end{equation}

\noindent where $c_1$ and $c_2$ are scaling coefficients. PPO is well suited to
the vehicle control domain because: (i) the continuous action space
(steer, throttle, brake) maps naturally onto the Gaussian policy it employs;
(ii) it is on-policy, avoiding the distribution-shift issues that complicate
off-policy methods in safety-critical settings; and (iii) its
data-collection strategy (rollout buffers of fixed length) integrates
cleanly with a simulation that runs at real time in a separate process.

The Stable-Baselines3 library \cite{Raffin2021} provides the PPO implementation
used in this project; it has been validated against reference implementations
across a range of benchmark tasks.

\subsection*{2.3 \quad SIMULATION-TO-REAL TRANSFER IN AUTONOMOUS DRIVING}
\addcontentsline{toc}{subsection}{2.3 \quad Simulation-to-Real Transfer in Autonomous Driving}

\indent Training autonomous driving policies in simulation rather than on physical
hardware is standard practice, motivated by safety, cost, and the ability to
reset episodes deterministically \cite{Kiran2021}. The primary risk is the
\emph{sim-to-real gap}: discrepancies between simulated and real sensor readings,
physics, and environmental conditions cause policies that perform well in
simulation to degrade when deployed.

Prominent simulation platforms include CARLA \cite{Dosovitskiy2017}, an
Unreal-Engine-based open-source simulator that models urban driving at high
visual fidelity, and the Unity ML-Agents Toolkit \cite{Juliani2020}, which
provides a general-purpose RL training framework tightly integrated with the
Unity Editor. SHILATE uses raw Unity physics without ML-Agents, for two
reasons: first, it avoids an opaque intermediate API layer and gives direct
control over the observation and reward structures; second, and more
importantly, the training process lives entirely in Python and communicates
through a standard MQTT broker, which is precisely the protocol used by the
target SDV deployment stack. The simulation therefore serves simultaneously as
a training environment and as an emulation of the vehicle bus interface.

\subsection*{2.4 \quad SOFTWARE-DEFINED VEHICLES AND THE VSS DATA MODEL}
\addcontentsline{toc}{subsection}{2.4 \quad Software-Defined Vehicles and the VSS Data Model}

\indent The COVESA Vehicle Signal Specification (VSS) \cite{COVESA2023} is a
hierarchical, vendor-neutral data model for vehicle signals, adopted by the
Eclipse SDV working group as the canonical representation for in-vehicle data.
VSS paths follow a dotted-name convention, e.g., \texttt{Vehicle.Speed} (km/h),
\texttt{Vehicle.Powertrain.CombustionEngine.Speed} (RPM), and
\texttt{Vehicle.Chassis.SteeringWheel.Angle} (degrees).

The Eclipse Kuksa databroker \cite{EclipseKuksa2024} implements a VSS-aligned
gRPC/REST server that acts as the vehicle signal hub. Producers (e.g., an
MQTT bridge receiving simulator telemetry) write to VSS paths; consumers (e.g.,
a safety application) subscribe to changes or poll current values. This
decoupled producer/consumer model is a cornerstone of the SDV architecture
used in this project.

Eclipse Velocitas \cite{EclipseVelocitas2024} builds on top of Kuksa by
providing a vehicle-application SDK and a deployment scaffold. A Velocitas app
subscribes to VSS signals, implements business logic, and publishes alerts or
control signals back to the vehicle bus. In SHILATE, the Velocitas application
monitors the live RL agent for speed and RPM limit violations.

\subsection*{2.5 \quad MQTT AND PUBLISH-SUBSCRIBE MESSAGING}
\addcontentsline{toc}{subsection}{2.5 \quad MQTT and Publish-Subscribe Messaging}

\indent MQTT 3.1.1 \cite{OASIS2014} is a lightweight publish-subscribe protocol
designed for constrained IoT environments. A central broker (Mosquitto, in
this project) decouples producers and consumers: a publisher pushes a message
to a topic string, and all subscribers to that topic receive it, with no
direct network coupling between the two parties. The protocol supports three
Quality-of-Service (QoS) levels; SHILATE uses QoS 0 (at-most-once delivery)
throughout, which minimises latency at the cost of possible message loss ---
an acceptable trade-off in a training loop where the next observation replaces
the current one within tens of milliseconds.

The MQTT protocol is also natively supported by the Eclipse Leda SDV image,
which ships Mosquitto as a core service. This means the same communication
fabric used during simulation training can be reused without modification on
the deployment target, eliminating the need for a protocol-translation layer
between training and production.

\subsection*{2.6 \quad RELATED WORK}
\addcontentsline{toc}{subsection}{2.6 \quad Related Work}

\indent A substantial body of work addresses RL for autonomous driving within
simulation \cite{Kiran2021}. OpenAI's \texttt{gym-car-racing} environment
\cite{Brockman2016} provides a top-down racetrack with pixel observations; it
is widely used as a benchmark but offers no path to real vehicle deployment.
CARLA-based RL studies \cite{Dosovitskiy2017} achieve high visual realism but
operate as closed systems with no integration to SDV middleware.

The Unity ML-Agents toolkit \cite{Juliani2020} enables RL training within
Unity but couples the Python trainer tightly to Unity's \texttt{side-channel}
API, making it difficult to reuse the same Python environment against a
physical vehicle. SHILATE deliberately avoids this coupling: the Python
environment communicates with Unity only through MQTT, the same protocol the
Leda SDV image uses for vehicle bus telemetry.

From the SDV side, the Eclipse Kuksa and Velocitas communities have
demonstrated containerised vehicle applications that read VSS signals and
enforce policies, but without an RL training pipeline attached. SHILATE
contributes the missing link: a training-to-deployment chain that is agnostic
to both the simulation engine and the SDV stack as long as both speak MQTT
and VSS.

\clearpage

% ============================================================
%  PLACEHOLDER: Chapters 3–8 will be added in subsequent stages
% ============================================================

\end{document}
